\sectionheading{5.31 Python Syntax-এর প্রাথমিক ধারণা}

Programming Language-এ সঠিকভাবে Code লেখার জন্য কিছু
নির্দিষ্ট নিয়ম অনুসরণ করতে হয়।

এই নিয়মগুলোর সমষ্টিকে \textbf{Syntax} বলা হয়।

%================================================
% Definition
%================================================

\booksubheading{Syntax কী?}

\begin{definitionbox}

\textbf{\color{darkTeal}সংজ্ঞা}

কোনো Programming Language-এ Statement বা Program
সঠিকভাবে লেখার জন্য নির্ধারিত নিয়মকে
\textbf{Syntax} বলে।

\end{definitionbox}

\textbf{সহজভাবে মনে রাখবে—}

\begin{center}

\textbf{Programming Language-এর Grammar}
=
\textbf{Syntax}

\end{center}

যেমন বাংলা বা English বাক্য লেখার নির্দিষ্ট নিয়ম আছে,
তেমনি Python Code লেখারও নির্দিষ্ট নিয়ম রয়েছে।

%================================================
% Statement
%================================================

\booksubheading{১. Python Statement}

Python Program-এর একটি কার্যকর Instruction-কে সাধারণভাবে
\textbf{Statement} বলা হয়।

উদাহরণ—

\begin{codebox}

\begin{english}
\ttfamily
print("Hello")\\
x = 10
\end{english}

\end{codebox}

এখানে প্রতিটি লাইন একটি Python Statement।

%================================================
% Indentation
%================================================

\booksubheading{২. Indentation}

Python Syntax-এর সবচেয়ে গুরুত্বপূর্ণ বৈশিষ্ট্যগুলোর একটি হলো
\textbf{Indentation}।

কোনো Statement-এর শুরুতে যে ফাঁকা জায়গা রাখা হয় তাকে
\textbf{Indentation} বলে।

\begin{examplebox}{উদাহরণ}

\begin{codebox}

\begin{english}
\ttfamily
if 10 > 5:\\
\hspace*{4mm}print("Yes")
\end{english}

\end{codebox}

এখানে
\textenglish{\ttfamily print("Yes")}
Statement-এর শুরুতে ফাঁকা জায়গা রাখা হয়েছে।
এটিই Indentation।

\end{examplebox}

Python-এ একই Code Block-এর Statementগুলো বোঝাতে
Indentation ব্যবহার করা হয়।

\begin{rememberbox}

\textbf{\color{deepOrange}সহজভাবে মনে রাখবে}

\begin{center}

\textbf{Indentation}
\(\longrightarrow\)
\textbf{Code Block নির্ধারণ করে}

\end{center}

\end{rememberbox}

%================================================
% 4 Spaces
%================================================

\booksubheading{সাধারণত কতটি Space ব্যবহার করা হয়?}

Python Code লেখার প্রচলিত Style অনুযায়ী একটি Indentation
Level-এর জন্য সাধারণত \textbf{4টি Space} ব্যবহার করা হয়।

\begin{codebox}

\begin{english}
\ttfamily
if 5 > 2:\\
\hspace*{4mm}print("Correct")
\end{english}

\end{codebox}

\begin{rememberbox}

\textbf{\color{deepOrange}সতর্কতা}

একই Block-এ Indentation অসামঞ্জস্যপূর্ণ হলে Error হতে পারে।

\textbf{ভুল উদাহরণ—}

\begin{codebox}

\begin{english}
\ttfamily
if 5 > 2:\\
print("Correct")
\end{english}

\end{codebox}

এখানে প্রয়োজনীয় Indentation নেই।

\end{rememberbox}

%================================================
% Colon
%================================================

\booksubheading{৩. Colon \textenglish{\ttfamily :}}

Python-এ
\textenglish{\ttfamily if},
\textenglish{\ttfamily for},
\textenglish{\ttfamily while},
\textenglish{\ttfamily def}
ইত্যাদির পরে নতুন Code Block শুরু করার ক্ষেত্রে সাধারণত
\textbf{Colon (\textenglish{\ttfamily :})}
ব্যবহার করা হয়।

\begin{codebox}

\begin{english}
\ttfamily
if x > 0:\\
\hspace*{4mm}print("Positive")
\end{english}

\end{codebox}

এখানে
\textenglish{\ttfamily if}
Condition-এর শেষে
\textenglish{\ttfamily :}
ব্যবহার করা হয়েছে।

\begin{rememberbox}

\textbf{\color{deepOrange}মনে রাখবে}

\begin{center}

\textbf{Block শুরু হওয়ার আগে}
\(\longrightarrow\)
\textenglish{\ttfamily :}
\(\longrightarrow\)
\textbf{Block-এর ভিতরে Indentation}

\end{center}

\end{rememberbox}

%================================================
% Case Sensitive
%================================================

\booksubheading{৪. Python Case Sensitive}

Python একটি
\textbf{Case Sensitive Programming Language}।

অর্থাৎ Capital Letter ও Small Letter-কে আলাদা হিসেবে
বিবেচনা করা হয়।

উদাহরণ—

\begin{codebox}

\begin{english}
\ttfamily
name\\
Name\\
NAME
\end{english}

\end{codebox}

উপরের তিনটি নাম Python-এর কাছে আলাদা।

একইভাবে—

\begin{codebox}

\begin{english}
\ttfamily
print("Hello")
\end{english}

\end{codebox}

সঠিক; কিন্তু—

\begin{codebox}

\begin{english}
\ttfamily
Print("Hello")
\end{english}

\end{codebox}

সাধারণ
\textenglish{\ttfamily print()}
Function হিসেবে কাজ করবে না।

\begin{rememberbox}

\textbf{\color{deepOrange}MCQ তথ্য}

Python
\(\longrightarrow\)
\textbf{Case Sensitive}

\end{rememberbox}

%================================================
% Quotation
%================================================

\booksubheading{৫. Quotation Mark}

Python-এ Text বা String লেখার জন্য সাধারণত
Quotation Mark ব্যবহার করা হয়।

\textbf{Double Quote-এর উদাহরণ—}

\begin{codebox}

\begin{english}
\ttfamily
print("Bangladesh")
\end{english}

\end{codebox}

\textbf{Single Quote-এর উদাহরণ—}

\begin{codebox}

\begin{english}
\ttfamily
print('Bangladesh')
\end{english}

\end{codebox}

অর্থাৎ Text লেখার ক্ষেত্রে সাধারণভাবে
Single Quote
\textenglish{\ttfamily ' '}
অথবা Double Quote
\textenglish{\ttfamily " "}
ব্যবহার করা যায়।

\begin{rememberbox}

\textbf{\color{deepOrange}সতর্কতা}

Quotation শুরু করলে সঠিকভাবে শেষ করতে হবে।

\textbf{ভুল—}

\begin{codebox}

\begin{english}
\ttfamily
print("Hello')
\end{english}

\end{codebox}

\textbf{সঠিক—}

\begin{codebox}

\begin{english}
\ttfamily
print("Hello")
\end{english}

\end{codebox}

\end{rememberbox}

%================================================
% Parentheses
%================================================

\booksubheading{৬. Parentheses \textenglish{\ttfamily ()}}

Python-এর অনেক Function ব্যবহার করার সময়
Parentheses
\textenglish{\ttfamily ()}
ব্যবহার করা হয়।

উদাহরণ—

\begin{codebox}

\begin{english}
\ttfamily
print("ICT")
\end{english}

\end{codebox}

এখানে—

\begin{itemize}

\item
\textenglish{\ttfamily print}
\(\longrightarrow\)
Function-এর নাম

\item
\textenglish{\ttfamily ()}
\(\longrightarrow\)
Parentheses

\item
\textenglish{\ttfamily "ICT"}
\(\longrightarrow\)
Function-এর ভিতরে দেওয়া Text

\end{itemize}

%================================================
% Comment
%================================================

\booksubheading{৭. Comment}

Python-এ এক লাইনের Comment সাধারণত
\textbf{\textenglish{\ttfamily \#}}
দিয়ে শুরু হয়।

\begin{codebox}

\begin{english}
\ttfamily
\# This is a comment\\
print("Hello")
\end{english}

\end{codebox}

\textenglish{\ttfamily \#}
-এর পরের অংশ Comment হিসেবে বিবেচিত হয় এবং Program-এর
কার্যকর Instruction হিসেবে Execute হয় না।

\begin{rememberbox}

\textbf{\color{deepOrange}মনে রাখবে}

\begin{center}

\textenglish{\ttfamily \#}
\(\longrightarrow\)
\textbf{Single-line Comment}

\end{center}

\end{rememberbox}

%================================================
% One Statement Per Line
%================================================

\booksubheading{৮. সাধারণত প্রতিটি Statement আলাদা লাইনে লেখা হয়}

Python Code পরিষ্কার ও সহজবোধ্য রাখার জন্য সাধারণত
প্রতিটি Statement আলাদা লাইনে লেখা হয়।

\begin{codebox}

\begin{english}
\ttfamily
a = 10\\
b = 20\\
print(a + b)
\end{english}

\end{codebox}

Python-এ কিছু ক্ষেত্রে Semicolon
\textenglish{\ttfamily ;}
ব্যবহার করে একই লাইনে একাধিক Statement লেখা সম্ভব।

উদাহরণ—

\begin{codebox}

\begin{english}
\ttfamily
a = 10; b = 20
\end{english}

\end{codebox}

তবে Beginner Program-এ পরিষ্কারভাবে আলাদা লাইনে লেখা ভালো—

\begin{codebox}

\begin{english}
\ttfamily
a = 10\\
b = 20
\end{english}

\end{codebox}

%================================================
% Code Block
%================================================

\booksubheading{৯. Code Block}

একই Control Structure বা Function-এর অন্তর্ভুক্ত
Statementগুলোর সমষ্টিকে সাধারণভাবে
\textbf{Code Block} বলা যায়।

উদাহরণ—

\begin{codebox}

\begin{english}
\ttfamily
if 10 > 5:\\
\hspace*{4mm}print("First")\\
\hspace*{4mm}print("Second")
\end{english}

\end{codebox}

এখানে Indentation করা দুটি
\textenglish{\ttfamily print()}
Statement একই Block-এর অংশ।

\begin{rememberbox}

\textbf{\color{deepOrange}সহজভাবে মনে রাখবে}

\begin{center}

\textenglish{\ttfamily :}
\(\longrightarrow\)
\textbf{Block শুরু}
\(\longrightarrow\)
\textbf{Indentation}
\(\longrightarrow\)
\textbf{Block-এর Statement}

\end{center}

\end{rememberbox}

%================================================
% Curly Braces
%================================================

\booksubheading{১০. Python-এ Code Block বোঝাতে Indentation ব্যবহৃত হয়}

C/C++-এর মতো অনেক Programming Language-এ Code Block
বোঝাতে Curly Braces
\textenglish{\ttfamily \{ \}}
ব্যবহার করা হয়।

উদাহরণস্বরূপ—

\begin{codebox}

\begin{english}
\ttfamily
\{\\
\hspace*{4mm}statements\\
\}
\end{english}

\end{codebox}

কিন্তু Python-এ Code Block বোঝাতে মূলত
\textbf{Indentation} ব্যবহার করা হয়।

\begin{rememberbox}

\textbf{\color{deepOrange}পরীক্ষায় গুরুত্বপূর্ণ}

\begin{center}

\textbf{C/C++}
\(\longrightarrow\)
\textenglish{\ttfamily \{ \}}

\qquad

\textbf{Python}
\(\longrightarrow\)
\textbf{Indentation}

\end{center}

\end{rememberbox}

%================================================
% Syntax Error
%================================================

\booksubheading{১১. Syntax Error}

Python-এর Syntax-এর নির্ধারিত নিয়ম ভুল হলে
\textbf{Syntax Error} হতে পারে।

\begin{examplebox}{Colon না দেওয়ার ভুল}

\textbf{ভুল—}

\begin{codebox}

\begin{english}
\ttfamily
if 10 > 5\\
\hspace*{4mm}print("Yes")
\end{english}

\end{codebox}

এখানে
\textenglish{\ttfamily if}
Statement-এর শেষে
\textenglish{\ttfamily :}
নেই।

\textbf{সঠিক—}

\begin{codebox}

\begin{english}
\ttfamily
if 10 > 5:\\
\hspace*{4mm}print("Yes")
\end{english}

\end{codebox}

\end{examplebox}

\begin{examplebox}{Parenthesis বন্ধ না করার ভুল}

\textbf{ভুল—}

\begin{codebox}

\begin{english}
\ttfamily
print("Hello"
\end{english}

\end{codebox}

\textbf{সঠিক—}

\begin{codebox}

\begin{english}
\ttfamily
print("Hello")
\end{english}

\end{codebox}

\end{examplebox}

%================================================
% Causes of Syntax Error
%================================================

\booksubheading{Syntax Error-এর সাধারণ কারণ}

\begin{compacttwocolpoints}

\item প্রয়োজনীয় Colon
\textenglish{\ttfamily :}
না দেওয়া।

\item প্রয়োজনীয় Indentation না দেওয়া।

\item Quotation Mark অসম্পূর্ণ রাখা।

\item Parenthesis
\textenglish{\ttfamily ()}
সঠিকভাবে বন্ধ না করা।

\item Keyword বা Function-এর বানান ভুল করা।

\item Capital ও Small Letter গুলিয়ে ফেলা।

\end{compacttwocolpoints}

%================================================
% Correct vs Incorrect
%================================================

\booksubheading{সঠিক ও ভুল Syntax-এর উদাহরণ}

\begin{tabularx}
{\linewidth}
{
|>{\RaggedRight\arraybackslash}p{0.29\linewidth}
|>{\RaggedRight\arraybackslash}p{0.29\linewidth}
|Y|
}

\hline

\rowcolor{softBlue}

\multicolumn{1}{|c|}{\textbf{ভুল}}
&
\multicolumn{1}{c|}{\textbf{সঠিক}}
&
\multicolumn{1}{c|}{\textbf{কারণ}}

\\

\hline

\textenglish{\ttfamily Print("Hi")}
&
\textenglish{\ttfamily print("Hi")}
&
Python Case Sensitive।
\\

\hline

\textenglish{\ttfamily if x > 5}
&
\textenglish{\ttfamily if x > 5:}
&
Colon অনুপস্থিত।
\\

\hline

\textenglish{\ttfamily print("Hi"}
&
\textenglish{\ttfamily print("Hi")}
&
Closing Parenthesis অনুপস্থিত।
\\

\hline

\textenglish{\ttfamily print("Hi')}
&
\textenglish{\ttfamily print("Hi")}
&
Quotation Mark মেলেনি।
\\

\hline

\textenglish{\ttfamily if x > 5:}\\
\textenglish{\ttfamily print(x)}
&
\textenglish{\ttfamily if x > 5:}\\
\hspace*{3mm}\textenglish{\ttfamily print(x)}
&
Block-এর Statement-এ Indentation প্রয়োজন।
\\

\hline

\end{tabularx}

%================================================
% Summary Table
%================================================

\booksubheading{Python Syntax এক নজরে}

\begin{tabularx}
{\linewidth}
{
|>{\bfseries\RaggedRight\arraybackslash}p{0.31\linewidth}
|Y|
}

\hline

\rowcolor{softBlue}

\multicolumn{1}{|c|}{\textbf{বিষয়}}
&
\multicolumn{1}{c|}{\textbf{নিয়ম}}

\\

\hline

Statement
&
সাধারণত আলাদা লাইনে লেখা হয়।
\\

\hline

Indentation
&
Code Block নির্ধারণ করতে ব্যবহৃত হয়।
\\

\hline

Colon
\textenglish{\ttfamily :}
&
Block শুরু করার আগে প্রয়োজনীয় স্থানে ব্যবহৃত হয়।
\\

\hline

Case
&
Python Case Sensitive।
\\

\hline

Text
&
Single বা Double Quotation Mark-এর মধ্যে লেখা যায়।
\\

\hline

Comment
&
\textenglish{\ttfamily \#}
দিয়ে Single-line Comment শুরু হয়।
\\

\hline

Function Call
&
সাধারণত Parentheses
\textenglish{\ttfamily ()}
ব্যবহৃত হয়।
\\

\hline

Code Block
&
Indentation দিয়ে প্রকাশ করা হয়।
\\

\hline

Syntax ভুল
&
Syntax Error হতে পারে।
\\

\hline

\end{tabularx}

%================================================
% Common Mistakes
%================================================

\booksubheading{সাধারণ ভুল}

\begin{compacttwocolpoints}

\item Python-এ Indentation শুধু Code সুন্দর করার জন্য মনে করা।

\item
\textenglish{\ttfamily if},
\textenglish{\ttfamily for},
\textenglish{\ttfamily while}
ইত্যাদির পরে প্রয়োজনীয়
\textenglish{\ttfamily :}
ভুলে যাওয়া।

\item
\textenglish{\ttfamily print()}
ও
\textenglish{\ttfamily Print()}
-কে একই মনে করা।

\item Single ও Double Quote অসম্পূর্ণভাবে ব্যবহার করা।

\item C/C++-এর মতো Python Code Block-এ
\textenglish{\ttfamily \{ \}}
ব্যবহার করা।

\item একই Block-এর Indentation অসামঞ্জস্যপূর্ণ করা।

\item
\textenglish{\ttfamily \#}
-কে Program Output-এর অংশ মনে করা।

\end{compacttwocolpoints}

%================================================
% Final Reminder
%================================================

\begin{rememberbox}

\textbf{\color{deepOrange}মনে রাখবে}

\begin{itemize}

\item
\textbf{Syntax}
হলো Code লেখার নিয়ম।

\item
Python-এ
\textbf{Indentation}
অত্যন্ত গুরুত্বপূর্ণ।

\item
একটি Indentation Level-এর জন্য সাধারণত
\textbf{4টি Space}
ব্যবহার করা হয়।

\item
\textenglish{\ttfamily if},
\textenglish{\ttfamily for},
\textenglish{\ttfamily while},
\textenglish{\ttfamily def}
ইত্যাদির Block শুরু করার আগে
\textenglish{\ttfamily :}
ব্যবহৃত হয়।

\item
Python
\textbf{Case Sensitive}।

\item
Text-এর জন্য Single বা Double Quotation Mark ব্যবহার করা যায়।

\item
\textenglish{\ttfamily \#}
দিয়ে Single-line Comment শুরু হয়।

\item
Python Code Block বোঝাতে
\textenglish{\ttfamily \{ \}}
-এর পরিবর্তে
\textbf{Indentation}
ব্যবহার করা হয়।

\item
Syntax-এর নিয়ম ভুল হলে
\textbf{Syntax Error}
হতে পারে।

\item
দ্রুত মনে রাখবে—

\begin{center}

\textbf{Syntax}
\(\longrightarrow\)
\textbf{Code লেখার নিয়ম}

\medskip

\textenglish{\ttfamily :}
\(\longrightarrow\)
\textbf{Block শুরু}
\(\longrightarrow\)
\textbf{Indentation}
\(\longrightarrow\)
\textbf{Block-এর Code}

\medskip

\textenglish{\ttfamily \#}
\(\longrightarrow\)
\textbf{Comment}

\qquad

\textbf{Python}
\(\longrightarrow\)
\textbf{Case Sensitive}

\medskip

\textbf{ভুল Syntax}
\(\longrightarrow\)
\textbf{Syntax Error}

\end{center}

\end{itemize}

\end{rememberbox}

\small

%================================================
% 5.32 Variable and Data Type
%================================================

